\section*{Chapter \ref{ch:rs:stylised-facts-of-returns} --- Stylised Facts of Returns}

\begin{solution}{exo:rs:stylised-facts-of-returns:1}
$-17.4\%/1.08\% = -16.2$ standard deviations; $\Phi(-16.2) \approx 6 \times 10^{-59}$, so about $10^{-54}$ such days in 26\,296: never, in
any number of universes' lifetimes. It happened once.
\end{solution}

\begin{solution}{exo:rs:stylised-facts-of-returns:2}
For independent summands the excess kurtosis divides by the number of terms: $3 + 17/5 = 6.4$, $3 + 17/21 = 3.8$, $3 + 17/63 = 3.3$.
The century's 11.5, 9.5 and 7.3 are far above: large moves cluster in time, so they do not average out.
\end{solution}

\begin{solution}{exo:rs:stylised-facts-of-returns:3}
$\P(X > x) \propto x^{-\alpha}$, so doubling the threshold divides the probability by $2^\alpha$: with $\alpha = 3$, one day in 8\,000 falls by
more than 8\%; with $\alpha = 4$, one in 16\,000.
\end{solution}

\begin{solution}{exo:rs:stylised-facts-of-returns:4}
For a Student-$t$ with $\nu$ degrees of freedom, $\E[T^2] = \nu/(\nu - 2)$ and $\E[T^4] = 3\nu^2/((\nu - 2)(\nu - 4))$, so the kurtosis is
$3(\nu - 2)/(\nu - 4) = 3 + 6/(\nu - 4)$, unchanged by scaling. A kurtosis of 19 needs $6/(\nu - 4) = 16$, $\nu = 4.375$: barely more than
four degrees of freedom.
\end{solution}

\begin{solution}{exo:rs:stylised-facts-of-returns:5}
The variance of an equal-weighted portfolio of 10 assets with volatility $\sigma$ and pairwise correlation $\rho$ is $\sigma^2(1/10 +
0.9\rho)$: $\sqrt{0.649} = 0.806\sigma$ on down days and $\sqrt{0.604} = 0.777\sigma$ on up days. Diversification removes 19\% of the
volatility when the market falls and 22\% when it rises.
\end{solution}

\begin{solution}{exo:rs:stylised-facts-of-returns:6}
Persistence $\alpha + \gamma/2 + \beta = 0.99$ (half the negative shocks carry $\gamma$); a variance shock halves in $\ln 0.5/\ln 0.99 = 69$
days. The autocorrelation of absolute returns decays like the persistence raised to the lag, exponentially: small after a few
half-lives, while the real series decays like a power of the lag.
\end{solution}

\begin{solution}{exo:rs:stylised-facts-of-returns:7}
With $\alpha = 0.01$, $\gamma = 0.05$, $\beta = 0.965$ the persistence is exactly one (integrated GARCH): the autocorrelations of $|m|$
at lags 20 and 100 rise to 0.37 and 0.26, against 0.21 and $-0.02$ for the default. But the variance no longer has a level to return
to (the intercept $\omega$ is zero): over the ten years the conditional volatility wanders between 1.7\% and 23.6\% a year and the
realised volatility is 7.7\% instead of the intended 16\%. Long memory needs another mechanism (a slow volatility component), not a
unit root.
\end{solution}

\begin{solution}{exo:rs:stylised-facts-of-returns:8}
A tail hedge pays in multi-day crashes and in sell-offs where correlations rise; a simulator calibrated to daily kurtosis and
clustering misses both (the book's reproduces daily tails but has monthly kurtosis 3.8 against 9.5 and no \omterm{def:rs:stylised-facts-of-returns:exceedance}{correlation asymmetry}). Its
tail events are single days, too short and too diversifiable, so the hedge's value is mismeasured; and ten-year windows of one
simulated path hold very few tail events. Validate on data with crashes, and on a simulator that reproduces the monthly tails.
\end{solution}

\begin{solution}{pb:rs:stylised-facts-of-returns:1}
\textbf{1.} 19.1 and 20.5. \textbf{2.} 2.90 and 2.83. \textbf{3.} 103, against 0.015. \textbf{4.} 22.0 (a fall of 17.4\% with a
prior standard deviation of 0.79\%). \textbf{5.} 8.7: a stock's specific shocks have Student-$t$ tails with four degrees of freedom
but only a part of its variance comes from the clustering market factor. \textbf{6.} Real 0.30, 0.22, 0.13; synthetic market 0.23,
0.21, $-0.02$. \textbf{7.} At lag 100: GARCH decays exponentially (half-life 69 days), the real series like a power. \textbf{8.} Yes:
$-0.111$ against the real $-0.085$. \textbf{9.} 9.5 against 3.8: the century's crashes lasted weeks (1929--1933, 1987, 2008) and
volatility regimes lasted years; the simulator has neither. \textbf{10.} Yes: 0.22 for $|r|$ against 0.11 for $r^2$.
\textbf{11.} 0.61 and 0.56. \textbf{12.} 0.20 on down days and 0.25 on up days: the simulator orders them the wrong way, so it
understates the cost of a crash to a diversified book. \textbf{13.} 44.5, 22\% of the total, against an edge of 3.58.
\textbf{14.} Four: the market and the strongest industry factors. \textbf{15.} Predictor measurement, backtest mechanics, risk models
at daily and weekly horizons, portfolio construction, costs. \textbf{16.} Tail hedging, drawdowns over months, diversification in
sell-offs, anything that depends on memory beyond fifty days. \textbf{17.} Multi-day crashes with rising correlations (downside
betas with a compensating drift, or a regime-switching volatility level). \textbf{18.} \emph{Named result:} the simulator matches the
real daily kurtosis (20.5 against 19.1), the left tail index (2.83 against 2.90) and the clustering at twenty days (0.21 against 0.22);
it misses the memory at a hundred days ($-0.02$ against 0.13), the monthly kurtosis (3.8 against 9.5) and the \omterm{def:rs:stylised-facts-of-returns:exceedance}{correlation asymmetry}
(0.20 down against 0.25 up, where the data give 0.61 against 0.56). \textbf{19.} Because a result is only as realistic as the facts
its data reproduce; the scorecard says which results transfer to markets. \textbf{20.} It is a test that any model of returns, or any
simulator, must pass before its conclusions are trusted.
\end{solution}

\begin{solution}{iq:rs:stylised-facts-of-returns:1}
Heavy tails (tail index about 3); no linear autocorrelation of returns but long-lasting autocorrelation of absolute returns
(volatility clustering, long memory); the leverage effect; \omterm{def:rs:stylised-facts-of-returns:aggregational}{aggregational Gaussianity} that is slow; \omterm{def:rs:stylised-facts-of-returns:gainloss}{gain--loss asymmetry}; correlations
that rise in sell-offs; \omterm{def:rs:stylised-facts-of-returns:intraday}{intraday seasonality}.
\iqlookfor{several facts, with a number or a mechanism for each.}
\end{solution}

\begin{solution}{iq:rs:stylised-facts-of-returns:2}
No: uncorrelated is not independent. Check the autocorrelations of $|r|$ and $r^2$ (volatility clustering), a test for ARCH effects,
and whether kurtosis falls with aggregation as fast as independence predicts.
\iqlookfor{the distinction, and a concrete test of nonlinear dependence.}
\end{solution}

\begin{solution}{iq:rs:stylised-facts-of-returns:3}
Correlations are higher on down days than on up days (0.61 against 0.56 among US industries at one standard deviation), and
volatility rises with them, so a diversified portfolio's volatility is higher exactly when losses happen; the common market factor
dominates in a sell-off.
\iqlookfor{\omterm{def:rs:stylised-facts-of-returns:exceedance}{correlation asymmetry}, not only rising volatility.}
\end{solution}

\begin{solution}{iq:rs:stylised-facts-of-returns:4}
Score it on the \omterm{def:rs:stylised-facts-of-returns:fact}{stylised facts} that matter to the strategy: tails, clustering and its memory, leverage, correlation structure and its
asymmetry, \omterm{def:rs:stylised-facts-of-returns:intraday}{intraday seasonality} if it trades intraday, against real data; then check that the strategy's result does not depend on a
fact the simulator misses (vary that feature and see if the answer moves).
\iqlookfor{a scorecard tied to the strategy's exposures.}
\end{solution}

\begin{solution}{iq:rs:stylised-facts-of-returns:5}
$3 + 17/21 = 3.8$ if independent. The observed 9 means that the daily returns are dependent in their sizes: large moves come in
clusters that last weeks, so monthly returns inherit heavy tails; volatility clustering, not only heavy daily tails.
\iqlookfor{the $1/n$ rule for excess kurtosis and the inference about dependence.}
\end{solution}

\begin{solution}{iq:rs:stylised-facts-of-returns:6}
The negative correlation between returns and later volatility: falls raise volatility more than rises. Measure the correlation of
$r_t$ with $r_{t+k}^2$ (or with realised volatility) for $k > 0$, against $k < 0$; or fit an asymmetric GARCH and test $\gamma$.
\iqlookfor{the direction of time in the measurement.}
\end{solution}
